\documentclass[11pt]{article}
\usepackage{amsmath}
\usepackage{amssymb}
\usepackage{mathpazo}
\usepackage[dvipsnames]{xcolor}
\usepackage[colorlinks, linkcolor=Blue]{hyperref}
\usepackage{graphicx}
\usepackage[margin=1in]{geometry}
\title{Homework 11}
\author{Your name here}
\date{\today{}}
\newcommand{\w}{\mathbf{w}}
\newcommand{\what}{\widehat{\w}}

\begin{document}
\maketitle

\section{Exact posterior for coin game}

Running the script \texttt{hw.py} will call the function \texttt{run\_scenario}
for each of the three prior scenarios we discussed in class, and this in turn
will generate for each scenario a plot of the prior and posterior densities for
$R$. First, fill in the missing pieces of code in \texttt{hw.py}.

Then, in your written solution, include these three plots; describe what the priors
represent, and explain the differences between the priors and posteriors (why do
they have the shapes they do). Also explain what makes the posteriors between
the three scenarios not the same.

%%% Answer START %%%
\textcolor{RoyalBlue}{Grading: 1 point for a reasonable explanation of what the
figures represent (as exemplified below).}
%%% Answer END %%%

\begin{figure}[htbp]
    \centering
    \includegraphics[width=0.3\textwidth]{Scenario_1.pdf}
    \includegraphics[width=0.3\textwidth]{Scenario_2.pdf}
    \includegraphics[width=0.3\textwidth]{Scenario_3.pdf}
%%% Answer START %%%
    \caption{\textcolor{RoyalBlue}{
        The three figures plot the prior and posterior densities for r under the three
        different scenarios, but each based on observations of 14 heads ($y_N$) out of 20
        coin tosses ($N$). (a) Scenario 1 starts with a uniform prior, so the
        observations have a big impact on the posterior belief, with majority of the
        density distributed around the ratio of 14 heads out of 20 (peaked at r = 7/10 =
        0.7); the variance is highest relative to the other two scenarios as the prior
        is so weak. (b) Scenario 2 has a fairly strong prior belief that the coin is
        completely fair, so it is peaked over r = 0.5 and the variance is fairly small;
        the evidence is that there are significantly more heads than tails, so the
        posterior is naturally pulled to the right, and the variance increases just
        slightly as the evidence is counter to the original prior belief, but not
        significantly. (c) Scenario 3 posits a relatively weak prior belief that the
        coin is biased towards heads. In light of the evidence, the posterior shifts to
        the right, close but slightly to the right of r = 0.7, and in general the
        variance has reduced.
        }
    }
    \label{fig:scenarios}
%%% Answer END %%%
\end{figure}

\newpage

\section{Workload}

How much time did you spend on this homework assignment outside of class?

\section*{Acknowledgments}

Cite all the people you've worked with on this homework, as well as any other
resources you used apart from the FCML textbook. If you used generative AI
tools (e.g., ChatGPT), please describe how you used them.
If you did not work with anyone else or use generative AI tools, please say so.

\end{document}
